\documentclass[../report.tex]{subfiles}
\begin{document}

\section{Số chu kỳ của ba phiên bản}

Bảng~\ref{tab:cycles} lấy từ golden model, đã được kiểm tra khớp từng chu kỳ với Logisim ở mục~\ref{sec:testres},
không tính đoạn 31 lệnh \texttt{sw} mà \file{run\_tests.py} chèn thêm. Cột ``NOP'' là số lệnh \texttt{nop} bản NH
thực sự thực thi; cột ``Load-use'' và ``Nhảy'' là số lần stall và số lần rẽ nhánh hoặc nhảy được thực hiện trên PL.

\begin{table}[!htb]
\centering\small
\caption{Số lệnh và số chu kỳ trên SC, NH và PL}\label{tab:cycles}
\setlength{\tabcolsep}{4.5pt}
\input{generated/cycles_table.tex}
\end{table}

Trên SC, CPI luôn bằng 1. Trên PL, CPI tính trên cả 12 chương trình là \resTotCPI. CPI cao nhất là \resMaxCPI{}
ở \resMaxProg: chương trình chỉ có \resMaxN{} lệnh nhưng nhảy \resMaxBR{} lần, nên \resMaxFlush{} chu kỳ flush
và 4 chu kỳ lấp đầy pipeline chiếm tỉ lệ lớn. \file{t09\_recursion.s}, chương trình dài nhất, có CPI
\resNineCPI. Bản NH có CPI tính trên số lệnh gốc là \resTotCPINH: nó thực thi \resTotNOP{} lệnh \texttt{nop},
trong khi PL chỉ stall \resTotLU{} chu kỳ cho load-use.

Chu kỳ thêm của PL so với số lệnh gồm 2 chu kỳ flush cho mỗi lần nhảy, 1 chu kỳ stall cho mỗi lần load-use và 4
chu kỳ lấp đầy pipeline. Trên \file{t09\_recursion.s}, \resNineBR{} lần nhảy tốn \resNineFlush{} chu kỳ, gấp hơn 5
lần \resNineLU{} chu kỳ của load-use: chương trình đệ quy có nhiều lệnh gọi, trả về và rẽ nhánh, nên control hazard
là nguồn mất mát lớn nhất. Ở NH, mọi phụ thuộc gần đều phải chờ bằng \texttt{nop}, kể cả phụ thuộc ALU sang ALU mà
PL forward được không mất chu kỳ nào. Với \file{demo.s} nạp sẵn trong ROM, cả ba bản kết thúc với TTY in
\texttt{Hello RISC-V!}, \texttt{a0} = \texttt{0xE8}, LED 7 đoạn hiện \texttt{00E8} và 8 LED sáng.

\section{Thời gian thực thi theo giả định trễ}

Chu kỳ clock của ba bản khác nhau, nên hình~\ref{fig:speedup} quy số chu kỳ ra thời gian theo giả định trễ ở
mục~\ref{sec:critical}: \resTSingle~ps cho SC, \resTPipe~ps cho NH và PL. Đây là ước tính theo giả định, không
phải số đo, vì Logisim không mô phỏng trễ cổng.

\begin{figure}[!htb]
\centering
\includegraphics[width=0.78\textwidth]{speedup.pdf}
\caption{Tăng tốc so với SC theo giả định trễ; đường đứt nét là mức lý tưởng khi CPI = 1}\label{fig:speedup}
\end{figure}

Tính trên tổng 12 chương trình, PL nhanh hơn SC \resSpeedPL{} lần và NH nhanh hơn SC \resSpeedNH{} lần, so với mức
lý tưởng \resIdealSpeedup{} lần. PL thực thi ít chu kỳ hơn NH \resSpeedPLNH{} lần với cùng chu kỳ clock; đây là
phần đóng góp riêng của \sig{ForwardUnit} và load-use stall bằng phần cứng. Khoảng cách tới mức lý tưởng có hai
nguyên nhân: CPI của PL lớn hơn 1 vì flush và stall, và năm tầng không dài bằng nhau (ID, WB chỉ cần 100~ps nhưng
vẫn chiếm trọn \resTPipe~ps). Đưa quyết định rẽ nhánh lên ID hoặc thêm dự đoán nhánh nhiều khả năng thu hẹp khoảng
cách này nhất trên các chương trình nhảy nhiều như \file{t09\_recursion.s}; mục~\ref{sec:future} ước lượng phần
lợi đó.

\end{document}
